\documentclass[11pt,a4paper]{article}
\usepackage[T1]{fontenc}
\usepackage[utf8]{inputenc}
\usepackage[french]{babel}
\frenchsetup{SmallCapsFigTabCaptions=false}
% The English language files of this Babel installation are broken, so the
% French language is loaded for hyphenation only and every caption is renamed.
\addto{\captionsfrench}{%
  \renewcommand{\contentsname}{Contents}%
  \renewcommand{\abstractname}{Abstract}%
  \renewcommand{\figurename}{Figure}%
  \renewcommand{\tablename}{Table}%
  \renewcommand{\refname}{References}%
}
\usepackage{lmodern,microtype,amsmath}
\usepackage[margin=23mm]{geometry}
\usepackage{booktabs,array,tabularx,longtable}
\usepackage{xcolor,tikz,enumitem,float}
\usepackage{hyperref}
\usetikzlibrary{arrows.meta,positioning,fit,calc,shapes.geometric}
\definecolor{ink}{HTML}{1F242C}
\definecolor{muted}{HTML}{5C626C}
\definecolor{line}{HTML}{D9DEE5}
\definecolor{accent}{HTML}{AF0032}
\definecolor{rose}{HTML}{FDF0F4}
\definecolor{agent}{HTML}{EDF1F8}
\definecolor{agentline}{HTML}{5A6B8C}
\definecolor{determ}{HTML}{F2F2F0}
\definecolor{warn}{HTML}{8A5A00}
\hypersetup{colorlinks=true,linkcolor=black,citecolor=black,urlcolor=ink,
 pdftitle={LANCE: where the project stands},
 pdfauthor={LANCE project},pdflang={en-GB}}
\urlstyle{same}
\def\UrlBreaks{\do\/\do\-\do\_\do\.\do\:}
\newcolumntype{Y}{>{\raggedright\arraybackslash}X}
\DeclareRobustCommand{\code}[1]{{\small\path{#1}}}
\newcommand{\source}[2]{\href{../../#1}{\nolinkurl{#1}}. #2}
% Engineering detail lives in the margin, so the argument stays readable.
\DeclareRobustCommand{\marginref}[1]{\marginpar{\raggedright\scriptsize\textcolor{muted}{#1}}}
\newcommand{\revision}{2026-11-02 17:00 (Europe/Brussels, UTC+01:00)}
\setlength{\parskip}{3pt}
\setlength{\emergencystretch}{2em}
\setlist{nosep,leftmargin=5mm,topsep=3pt}
\clubpenalty=10000 \widowpenalty=10000
\setcounter{tocdepth}{1}
% A clear boundary between the state (which is the point) and the evidence.
\usepackage{titlesec}
\titleformat{\section}[block]
  {\normalfont\sffamily\bfseries\Large\color{ink}}
  {\colorbox{ink}{\makebox[1.1em][c]{\color{white}\thesection}}}
  {0.8em}{}[{\vspace{3pt}{\color{line}\titlerule[0.6pt]}}]
\titlespacing*{\section}{0pt}{18pt}{12pt}

\title{LANCE\\[4pt]
\large Where the project stands\\[2pt]
\normalsize State report as of 2 November 2026}
\author{LANCE project\\\small Written by specialist strands, independently reviewed and synthesised}
\date{Last revision: \revision\\[3pt]
\small Documentation state as of 2 November 2026; no new experiment is reported}

\begin{document}
\maketitle

\begin{abstract}
LANCE is a platform for the security audit of authorised connected-object
networks, developed in the context of the UMONS work of the NATO SPS G6259
programme. It asks one measurable question: \emph{when does a language model find
something that a competent automation cannot find, given the same tools, the same
evidence and the same budgets?}

This report answers a reading question: where does the project stand today? Its
first pages give the answer as a whole --- what the system does, where the model
actually decides, which three policies are compared, and what blocks the
conclusion (Section~\ref{sec:etat}). The parts that follow are the evidence behind
that answer: the pipeline architecture, the audit protocols and their evidence
contracts, the experimentation harnesses and their sandbox, the execution
infrastructure and the delivery pipeline, and the benchmark itself.

The position is unchanged since the previous report of 30 October 2026: the system
is built, instrumented and documented, and its chain of evidence is sound. The
scientific question remains open. No superiority of the model is established, no
confidence interval can be computed, and the obstacle is no longer the
computation but the absence of independent families to compute on.

The document consistently distinguishes what is implemented, what is measured and
what remains proposed. No new experiment is reported: the measurements cited are
those already recorded, with their date. The limits of each component are
explicit, including those of the sandbox, whose guard is not an exhaustive
native-code sandbox.
\end{abstract}

\tableofcontents

% === WHERE THE PROJECT STANDS: the whole answer, before the evidence ===
\section{Where the project stands}
\label{sec:etat}

This section answers the question of the report before the evidence arrives. It
holds in six pages: what an audit does and who decides at each step
(Figure~\ref{fig:audit}), which three policies are compared and what is missing for
the comparison to be fair (Figure~\ref{fig:policies}), what blocks the conclusion
and how the project got there (Figure~\ref{fig:blocages} and
Table~\ref{tab:etat-chiffre}). The five parts that follow are the proof.

\subsection{The project in one paragraph}

A security auditor takes a network it has never seen, finds what is wrong with it,
proves each finding rather than guessing it, and writes a report an operator can
act on. LANCE does this with a language model in the loop. The programme asks a
precise question: \emph{the model earns its place only where a competent script
would fail --- does that situation exist, and can it be measured?}

The answer requires three things, and the report follows that order. A system
whose behaviour is traceable, so that a result can be attributed to a cause
(Figure~\ref{fig:audit}). A comparison in which the only thing that changes is the
reasoning capacity, not the information or the tools
(Figure~\ref{fig:policies}). Scenarios where the two candidates genuinely differ
(Figure~\ref{fig:blocages}). The first exists, the second is half built, the third
does not yet exist. That is the whole state of the project.

% ---------------------------------------------------------------------------
\subsection{What an audit does, and where the model decides}
\label{sec:parcours}

An audit runs in six phases. Each phase produces one document that the next one
reads, and nothing else crosses between them. Figure~\ref{fig:audit} is the map to
keep in mind for the rest of the report: it shows the six phases, what each
produces, and --- the point that matters --- \emph{who} decides at each step.

\begin{figure}[H]
\centering
\begin{tikzpicture}[
  x=1mm, y=1mm,
  phase/.style={draw=line, fill=white, line width=.55pt, rounded corners=2pt,
    text=ink, align=center, font=\sffamily\footnotesize, inner sep=4pt,
    minimum height=15mm, text width=22mm},
  model/.style={phase, draw=agentline, fill=agent, line width=.9pt},
  fixed/.style={phase, fill=determ},
  prod/.style={font=\sffamily\scriptsize, text=muted, align=center},
  flow/.style={-{Latex[length=1.8mm,width=1.0mm]}, draw=muted, line width=.6pt},
  lbl/.style={font=\sffamily\scriptsize, text=muted}
]
% row of phases
\node[fixed] (p1) at (0,0)   {\textbf{1. Graph}\\{\tiny\normalfont map the network}};
\node[model] (p2) at (28,0)  {\textbf{2. Recon}\\{\tiny\normalfont probe devices}};
\node[model] (p3) at (56,0)  {\textbf{3. Hypotheses}\\{\tiny\normalfont what could break}};
\node[model] (p4) at (84,0)  {\textbf{4. Verification}\\{\tiny\normalfont test it}};
\node[fixed] (p5) at (112,0) {\textbf{5. Intrusion}\\{\tiny\normalfont what is reachable}};
\node[fixed] (p6) at (140,0) {\textbf{6. Report}\\{\tiny\normalfont for the operator}};
% deliverables
\node[prod] (d1) at (0,-13)  {graph\\analysis};
\node[prod] (d2) at (28,-13) {recon log};
\node[prod] (d3) at (56,-13) {vulnerabilities};
\node[prod] (d4) at (84,-13) {exploit proofs};
\node[prod] (d5) at (112,-13){access map};
\node[prod] (d6) at (140,-13){final report};

\draw[flow] (p1) -- (p2);
\draw[flow] (p2) -- (p3);
\draw[flow] (p3) -- (p4);
\draw[flow] (p4) -- (p5);
\draw[flow] (p5) -- (p6);
\foreach \p/\d in {p1/d1,p2/d2,p3/d3,p4/d4,p5/d5,p6/d6}
  \draw[draw=muted!60, line width=.4pt, -{Latex[length=1.2mm,width=.9mm]}] (\p) -- (\d);

% the band that says who decides
\node[lbl, anchor=east] at (-13,2.5)   {who decides};
\draw[muted, line width=.4pt] (-13,5) -- (152,5);
\draw[fill=rose, draw=accent, line width=.6pt, rounded corners=1pt]
  (-11,12) rectangle (8,16.5);
\draw[fill=agent, draw=agentline, line width=.8pt, rounded corners=1pt]
  (14,12) rectangle (99,16.5);
\draw[fill=determ, draw=muted, line width=.6pt, rounded corners=1pt]
  (105,12) rectangle (151,16.5);
\node[font=\sffamily\tiny, align=center, text=accent] at (-1.5,14.25)
  {the code:\\rendered from the log};
\node[font=\sffamily\tiny, align=center, text=agentline] at (56.5,14.25)
  {\textbf{the model reasons here on evidence the tools collected}};
\node[font=\sffamily\tiny, align=center, text=muted] at (128,14.25)
  {the code:\\no free choice};

% legend
\node[font=\sffamily\tiny, anchor=west, text=muted] at (-13,-21)
  {\tikz[baseline=-0.6ex]{\fill[agent] (0,0) rectangle (.30,.12);\draw[agentline, line width=.8pt] (0,0) rectangle (.30,.12);}\ model decides};
\node[font=\sffamily\tiny, anchor=west, text=muted] at (4,-21)
  {\tikz[baseline=-0.6ex]{\fill[determ] (0,0) rectangle (.30,.12);\draw[muted, line width=.5pt] (0,0) rectangle (.30,.12);}\ code decides};
\node[font=\sffamily\tiny, anchor=west, text=muted] at (21,-21)
  {every phase writes its deliverable; the model's free text is never the source of a finding};
\end{tikzpicture}
\caption{An audit in six phases, and the boundary of the model's freedom. Only
phases 2, 3 and 4 leave a real choice to the model: which probes to run, which
weaknesses to hypothesise, and how to test each one. Phase 1 renders its result
from the graph tools, phase 5 from the tool log, and phase 6 assembles the report
deterministically from per-section sheets --- its agent declares no tool and runs a
single turn. The check that produces a score is a separate program: an offline
evaluator that recounts the run from its log and from an oracle fixed before the
run.}
\label{fig:audit}
\end{figure}

\textbf{The log is the source, not the model's conclusion.} Every tool call is
archived with a reference, a refusal or an exception included. A finding that
cites a piece of evidence which does not appear in the log is rejected. This is
what makes a result attributable: the report states what happened, not what the
policy believed happened.

\textbf{The oracle precedes the run.} The expected answer is fixed before the
policy runs, in a file the policy never receives. A finding that contradicts the
oracle invalidates the trial instead of correcting it. This precedence is what
prevents the answer from being adjusted after the fact.

\textbf{The score is recomputed, never taken on trust.} The evaluator receives
only the briefing, the oracle, the log and the final decision. It refuses to award
a success to a policy that lists a finding without the tool call that supports it.

\marginref{Where this lives: phase registry \code{src/agent/phases/registry.py}; governance boundary \code{src/agent/core/executor.py}; evaluator under \code{src/benchmark/}.}

% ---------------------------------------------------------------------------
\subsection{The three policies compared on the same scenario}
\label{sec:politiques}

The programme compares three ways of answering the same question on the same
scenario. They receive the same targets and are judged by the same evaluator; what
differs is how the answer is produced.

\begin{figure}[H]
\centering
\begin{tikzpicture}[
  x=1mm, y=1mm,
  col/.style={draw=line, fill=white, line width=.55pt, rounded corners=3pt,
    text=ink, align=center, font=\sffamily\footnotesize, inner sep=5pt,
    text width=38mm, minimum height=14mm},
  known/.style={col, draw=muted, fill=determ},
  missing/.style={col, draw=accent, line width=.9pt, fill=rose, dashed},
  head/.style={font=\sffamily\small\bfseries, text=ink, align=center},
  sub/.style={font=\sffamily\scriptsize, text=muted, align=center},
  arrow/.style={-{Latex[length=2mm,width=1.3mm]}, draw=muted, line width=.7pt},
  badge/.style={font=\sffamily\scriptsize, rounded corners=1.5pt, inner sep=2pt,
    align=center}
]
% top row: what every policy receives
\node[col, text width=44mm] (scen) at (0,42) {\textbf{the same scenario}\\{\tiny\normalfont deployed, reset, audited once}};
\node[col, text width=44mm] (obs) at (-52,42) {\textbf{the same observations}\\{\tiny\normalfont what the scanner collected\\ on each device}};
\node[col, text width=44mm] (eval) at (52,42) {\textbf{the same evaluator}\\{\tiny\normalfont oracle fixed before,\\ score recomputed}};

% every input flows down to all three policies
\coordinate (fanL) at (36,30);
\coordinate (fanR) at (72,30);
\node[font=\sffamily\tiny, text=accent, anchor=west] at (30,30.5) {each policy receives all three, and only these};
\foreach \p in {scen,obs,eval} {
  \draw[draw=muted!70, line width=.5pt] (\p.south) -- (\p.south |- fanL);
}
\draw[draw=muted!70, line width=.5pt] (scen.south |- fanL) -- (eval.south |- fanL);
\draw[arrow, accent, line width=.7pt] (-46,30) -- (-46,22);
\draw[arrow, accent, line width=.7pt] (0,30) -- (0,22);
\draw[arrow, accent, line width=.7pt] (46,30) -- (46,22);

% three policies
\node[known] (scripted) at (-46,10) {\textbf{\sffamily scripted}\\{\scriptsize\normalfont deterministic procedures:\\\normalfont discover, inventory, audit}};
\node[known] (rules) at (0,10)      {\textbf{\sffamily rules}\\{\scriptsize\normalfont a declared baseline:\\\normalfont scanner, then interpretation rules}};
\node[missing] (llm) at (46,10)     {\textbf{\sffamily llm\_on\_evidence}\\{\scriptsize\normalfont the model interprets the very\\\normalfont observations the rules receive}};

\node[badge, fill=determ, text=muted] at (-46,-3) {runs; no benchmark score};
\node[badge, fill=determ, text=muted] at (0,-3)   {runs; strict-v3 score};
\node[badge, fill=rose, text=accent] at (46,-3)   {\textbf{does not exist}};

\node[sub, text width=40mm] at (-46,-14) {no model at all};
\node[sub, text width=40mm] at (0,-14)
  {\textbf{the question}: same data as \texttt{rules},\\ a model where the rules are};
\node[sub, text width=46mm] at (46,-20)
  {\raggedright the third policy is what the whole programme needs; without it the contrast measures verification, not interpretation};
\end{tikzpicture}
\caption{The three policies and what is missing. \code{scripted} and \code{rules}
exist and are scored; \code{rules} is the declared baseline. The third policy,
\code{llm\_on\_evidence}, would give the model the observations the rules already
receive and nothing else --- no rule findings, no ground truth, no expected count
of flaws. It does not exist, and without it the comparison measures verification
rather than interpretation.}
\label{fig:policies}
\end{figure}

\textbf{Why the third policy is the one that matters.} An exploration where the
model roams a network freely (\code{vanilla}) changes too much at once: the
targets chosen, the order of the probes and the interpretation all differ, so a
better score cannot be attributed to reasoning. A policy that receives the
scanner's observations and must produce findings backed by a citation of those
observations isolates exactly one variable: \emph{who interprets}. That is the
question the programme asked.

\textbf{What currently favours the rules.} Four mechanisms give the declared
baseline an advantage the model does not have: the phase-four probes only the
rules may emit; a supplement that rewrites the scan files before aggregation; the
scanner findings added to every aggregation regardless of the policy; and a
benchmark score reserved for the full pipeline. Until one of these is removed or
mirrored, an equal score would say nothing and a lower score would be
uninterpretable.

% ---------------------------------------------------------------------------
\subsection{What blocks the conclusion}
\label{sec:blocages}

The obstacle is not the measurement machinery, which is now sound, but the
material it would need. Figure~\ref{fig:blocages} shows the chain as it stands and
the one observation that changed everything on 30 October 2026.

\begin{figure}[H]
\centering
\begin{tikzpicture}[
  x=1mm, y=1mm,
  pstep/.style={draw=line, fill=white, line width=.55pt, rounded corners=2pt,
    text=ink, align=center, font=\sffamily\footnotesize, inner sep=5pt,
    text width=31mm, minimum height=15mm},
  ok/.style={pstep, draw=muted, fill=determ},
  blocked/.style={pstep, draw=accent, line width=.9pt, fill=rose},
  later/.style={pstep, dashed, text=muted},
  flow/.style={-{Latex[length=2mm,width=1.2mm]}, draw=muted, line width=.7pt},
  ret/.style={-{Latex[length=1.8mm,width=1.0mm]}, draw=accent, line width=.7pt, dashed},
  note/.style={font=\sffamily\tiny, text=muted, align=left, anchor=north west}
]
% the chain, left to right
\node[ok] (f) at (0,20) {\textbf{a scenario}\\[-2pt]{\scriptsize\normalfont a deployment,\\ resettable}};
\node[ok] (h) at (48,20) {\textbf{a family}\\[-2pt]{\scriptsize\normalfont variants of one competence,\\ counted as one observation}};
\node[blocked] (c) at (96,20) {\textbf{a closure test}\\[-2pt]{\scriptsize\normalfont does a one-line rule\\ already solve it?}};
\node[later] (m) at (144,20) {\textbf{a measure}\\[-2pt]{\scriptsize\normalfont arm against arm,\\ interval, decision}};
\node[note, text width=34mm] at (117,11) {a family is built only if no generic rule crosses it};

\draw[flow] (f) -- (h);
\draw[flow] (h) -- (c);
\draw[flow, accent] (c) -- node[above, font=\sffamily\tiny, text=accent] {if it passes} (m);

% what happened on 30 October, below the closure test
\node[ok, text width=28mm, minimum height=11mm] (build) at (18,-8)
  {\scriptsize\normalfont built: a probe of\\ three registers};
\node[ok, text width=26mm, minimum height=11mm] (review) at (52,-8)
  {\scriptsize\normalfont six reviews\\ validated it};
\node[blocked, text width=30mm, minimum height=11mm] (reject) at (89,-8)
  {\scriptsize\normalfont one hour of execution\\ rejected it};
\draw[flow] (build) -- (review);
\draw[flow] (review) -- (reject);
\draw[ret] (reject.north) -- node[right, font=\sffamily\tiny, text=accent, align=left,
  inner sep=2pt] {the trap was crossed by a rule \\ \(\Rightarrow\) nothing to measure} (c.south);

% what the episode says, in the margin of the chain
\node[align=left, anchor=north, text width=70mm, font=\sffamily\scriptsize, text=ink] at (26,-21)
  {\textbf{The one hour that mattered.} {\tiny\normalfont\itshape
   The family \texttt{s\_registres\_constructeur} asked whether a device commands
   something or merely caches a value. Its trap was meant to defeat the rules. A rule
   that reads the effect of a write on the neighbouring registers clears it three
   times out of three --- and follows from the public description of the device.}};
\end{tikzpicture}
\caption{The chain from a scenario to a measurement, and where it breaks. Every
stage exists except the measure; what is missing is families that a generic rule
does not cross. The episode of 30 October shows why the closure test moved ahead of
the reviews: the protocol rejects in an hour what six readings had validated.}
\label{fig:blocages}
\end{figure}

\textbf{What the closure test is.} Before building a family, one asks whether a rule
that solves the case can be written in a single line. If it can, the family measures
the omission of a declaration rather than a limit of interpretation, and it is not
built. The test is cheap; running it before writing eight dossiers is the change of
method this month produced.

\textbf{Why this is a blockage and not a delay.} The inference tool was repaired on
30 October: it now computes a one-sided p-value against the pre-registered
threshold, refuses to emit a figure below three families, and requires five to
decide. It is therefore correct, and idle: there is nothing to compute on. Building
more families of the same kind would not change the arithmetic.

% ---------------------------------------------------------------------------
\subsection{The state in one table}
\label{sec:table-etat}

\begin{table}[H]
\centering
\caption{The state of the project on 2 November 2026, and where the proof is.}
\label{tab:etat-chiffre}
\small
\begin{tabularx}{\linewidth}{@{}p{3.9cm}Yp{2.6cm}@{}}
\toprule
Question & Answer & \textbf{Proof} \\
\midrule
Is an audit traceable? & Yes. Every tool call is archived with a reference,
bounded, and a lost piece of evidence stops the run instead of leaving a silent
gap. The guards are application-level. & \textbf{Part~\ref{sec:architecture-pipeline}} \\
Does the harness run outside the laboratory? & Yes, without broker and without a
model, under a guard whose blind spots are declared. It instruments a contrast; it
measures nothing. & \textbf{Part~\ref{sec:harnais-sandbox}} \\
Are the protocols' contracts sound? & Yes. Briefing, private oracle and offline
evaluator are separated, and the oracle is fixed before the run. & \textbf{Part~\ref{sec:protocoles}} \\
Can results be delivered reproducibly? & Yes. Push, validate, deploy to one of
three instances; no execution on a workstation. Single points of failure declared.
& \textbf{Part~\ref{sec:infrastructure}} \\
Does the model beat the rules? & \textbf{Unknown, and not yet measurable.} & \textbf{Part~\ref{sec:benchmark}} \\
How much of the corpus needs interpretation? & 6 entries out of 288 (2.1\,\%)
by the mechanical reading; 84.4\,\% detectable by signature by the interpretive
reading. Two counts, never to be cited one for the other. & \textbf{Part~\ref{sec:benchmark}} \\
Is the measurement statistically valid? & The tool is repaired and refuses an
undefensible figure; the families that would give it data do not exist. & \textbf{Part~\ref{sec:benchmark}} \\
\bottomrule
\end{tabularx}
\end{table}

\subsection{What to read next}

A reader who needs the whole story in ten minutes reads this section only, then
the conclusion, which states what the project can and cannot assert. A reader who
wants to check a claim follows the cross-reference in
Table~\ref{tab:etat-chiffre}. Each of the five parts opens with two paragraphs
that summarise it, and keeps the engineering detail --- file names, functions,
commit identifiers --- in its margins, so that the argument can be read without
them.

\clearpage


% === VOCABULARY (early, so no later section introduces a term) ===
\section{Vocabulary}

The report uses a small vocabulary throughout. Each term is defined here, in the
order a newcomer meets them, so that no later section has to introduce it.

\subsection{What is compared}

\begin{description}[leftmargin=3.1cm, style=nextline]
\item[Policy (politique)] One way of carrying out an audit. The programme
compares three: a declared automation, a second automation that follows
convention, and an agent driven by a language model.

\item[Arm (bras)] One concrete implementation of a task inside a family. The
arms share the briefing and the journal; only the interpreter changes. A family
is typically run with three: the declared rules, a script reading by convention
and the model. The witness is a value inside the fixture, not an arm.

\item[D1] The declared automation: the rules written explicitly for the task.
\item[D1+] D1 \emph{augmented with every rule derivable from the public inputs}.
This is the reference a model must beat for a result to mean anything. If a
competent reader who knows the briefing can state the rule, D1+ contains it, and
the family measures nothing.

\item[Family (famille)] A group of variants testing the same competence. Its
variants count as a single observation, because they test one thing. Two
families share no simulator, no vocabulary and no failure mechanism.

\item[Closure test (test de clôture)] The check performed before measuring a
family: a competent automation must fail. If it succeeds, the family measures
nothing about the model and is rejected. The test must be executable code, not a
doctrine.

\item[Filter (filtre préalable)] The cheaper check that precedes construction: if
the rule solving the case can be written in one line, a script will have it.
\end{description}

\subsection{What makes a result verifiable}

A result is verifiable when it can be recounted by someone who does not trust the
system. The report uses five objects for that, each answering one question.

\begin{description}[labelwidth=22mm, leftmargin=26mm, style=nextline, labelsep=2mm]
\item[Briefing (public)] Everything given to an arm before it runs: the task, the
target, the tools. It never contains the posture or the expected verdict.
\item[Oracle (private)] The expected answer, written before the run and never shown
to the arms. It states which state was installed and which fact is observable, never
which decision the policy should return.
\item[Witness (témoin)] A fresh value drawn when the fixture starts, whose value
appears nowhere in the briefing. It proves that what was observed is fresh and not a
stale artefact.
\item[Journal (journal)] The archived log of every tool call made during a run, with
its arguments, result and reference. It is what the evaluator recounts; the policy's
own conclusion is not trusted.
\item[Evaluator (offline)] The program that rereads the journal and recomputes the
outcome. A cited piece of evidence absent from the journal is rejected.
\end{description}

A \textbf{fixture} is the state installed on the target before a trial: a healthy
posture, a vulnerable one, or a frozen corpus read offline. A \textbf{policy} is one
way of carrying out an audit (a declared automation, an agent driven by a model). An
\textbf{arm} is one concrete implementation of a task inside a family; arms share
the briefing and the journal, and only the interpreter changes.

\subsection{How results are judged}

\begin{description}[leftmargin=3.1cm, style=nextline]
\item[Funnel] The three stages of a finding, in order: \texttt{candidates},
\texttt{filtered}, \texttt{confirmed}. Only the last counts as a true positive.

\item[Strict-v3] The normalised vulnerability contract and the metrics derived
from it. The \texttt{confirmed.*} metrics count only findings whose evidence has
been independently validated against the tool log.

\item[Campaign admission] The set of conditions a campaign must meet before its
results can be aggregated. When one is unsatisfied, the results remain
descriptive and are declared as such.

\item[V0] The validity gate placed before every hypothesis: a control task that
is solvable without inference. If either arm fails it, no downstream result is
readable.

\item[H1 to H6] The six hypotheses of the register, each naming a kind of
competence a model might have where a script does not. H2 is the only one
measured; H4 is the one the rejected family attempted.

\item[C2] A candidate scenario rejected after two adversarial reviews, on paper,
before any run. Its rejection is a published negative result.

\item[Posture] The state of a fixture with respect to security: healthy (no flaw
installed) or vulnerable (flaw installed).
\end{description}


\section{Scope and method}
\label{sec:scope}

This report answers a reading question: \emph{how is the system built, and where
does it actually stand?} It does not replace the reference pages of \code{docs/},
which describe the current behaviour, nor the detailed protocols of
\code{research/}, which carry the reasoning. It succeeds the report of
30 October 2026 and does not revise it: that folder stays as the archive of its
state.

The body was written by specialist strands working on the \code{dev/1} branch, then
reviewed against that same code by an independent reader. This revision changes
the order and the register: the state comes first and the evidence after, and the
engineering detail was moved to the margins so that the argument stays readable.
No figure of this report is new; all of them are recomputed from the repository or
reported as read in a document.

The report is dated: it describes a state, not a contract. Later behaviour may
contradict it without making it false. Statements that cannot be verified in the
repository are flagged as such rather than filled in.

% === EVIDENCE ===
\section{How an audit is built}
\label{sec:architecture-pipeline}

\textbf{In brief.} An audit goes through six phases, each writing one document
that the next one reads. The application is the thinnest part: it accepts a
request, allows one run per instance, streams events and reads the artefacts back.
The engine is a pipeline that wraps every tool call in a governance boundary: a
call is checked against the budgets, archived with a reference, and archived even
when it is refused or fails. The policy only ever reasons over evidence that the
tools produced, and the score comes from a separate evaluator, not from the policy.

\textbf{What the code enforces, and what it does not.} What is enforced: a lost
piece of evidence stops the run; a deliverable that misses a device is rejected;
the report is assembled deterministically; the phase-five status can never be
certified. What is not guaranteed: the budgets only exist when they are passed
explicitly, the laboratory lock is optional, and two scenarios on a shared
laboratory are not mutually excluded.

\marginref{\code{docs/architecture/run-modes.md} for the modes; \code{src/agent/run\_modes.py} for the selection.}

\subsection{The application does not carry the logic}

The web layer validates the request, guarantees a single active run per instance,
streams the events and re-reads the artefacts. The strategy is chosen by
\code{make\_runner}, not by the interface: for \code{lance} and \code{rules} it
instantiates a pipeline and wraps it in an adapter that reuses the same object and
its public API; for \code{vanilla} and \code{scripted} it builds a standalone
runner that shares with the pipeline only the tool boundary, the evidence log, the
cost tracking and the laboratory reservation.

\begin{table}[H]
\centering
\caption{Execution modes and their boundary.}
\label{tab:run-modes}
\small
\begin{tabular}{@{}llll@{}}
\toprule
Mode & Engine & Model & Effects \\
\midrule
\code{lance} & full pipeline, six phases & required & scenario, deployment, evaluation \\
\code{rules} & pipeline, phases 3--4 & none & scenario and deterministic audit \\
\code{vanilla} & standalone runner & required & bounded network probe on an address range \\
\code{scripted} & standalone runner & none & versioned procedures \\
\bottomrule
\end{tabular}
\end{table}

The \code{rules} mode is not a second engine: the factory sets the decision policy,
restricts the phases to analysis and verification, and refuses any provider, any
direct network access and any monetary budget. The two standalone modes, in turn,
refuse a scenario and every harness-specific setting. The qualifier
\code{full}/\code{compact} designates the orchestration, never the type of runner.

\subsection{The six phases and what each hands over}

The registry assigns to each phase a prompt, a tool group, a validator and a
deliverable; the routing is an explicit table, not a dynamic import.

\begin{table}[H]
\centering
\caption{Phases, main deliverables and associated artefacts.}
\label{tab:phases}
\small
\begin{tabular}{@{}cllp{4.4cm}@{}}
\toprule
Phase & Main deliverable & Who decides & Associated artefacts \\
\midrule
1 & \code{01\_graph\_analysis.md} & code (projection of the graph tools) & \code{01\_graph\_evidence.json} \\
2 & \code{02\_recon.md} & the model chooses the probes & \code{02\_recon\_evidence.json}, \code{02\_topology\_edges.json} \\
3 & \code{03\_vuln\_analysis.json} & the model writes the hypotheses & \code{03\_vuln\_analysis\_raw.json}, \code{03\_phase3\_status.json}, per-device sheets \\
4 & \code{04\_exploitation.json} & the model tests one flaw per micro-agent & \code{execution\_errors} diagnostics \\
5 & \code{05\_intrusion.json} & code (assessment rendered from the log) & \code{05\_intrusion\_context.json} \\
6 & \code{06\_report.md} & code (assembled from section sheets) & \code{06\_report\_analysis.md} \\
\bottomrule
\end{tabular}
\end{table}

Phase 1 projects the results of the graph tools into a log of facts independent of
the model's discourse, and the runner rejects a submission whose per-device
coverage is incomplete. Phase 2 lets the model choose the probes, but the recon
document is rendered from the tool log rather than copied by the model. Phase 3
runs in three steps: deterministic scanning, one analyst per device, then an
aggregation that normalises, deduplicates and projects the findings. Phase 4
launches one micro-agent per vulnerability. Phase 5 produces its assessment from
the log: the model supplies an end marker only, and the status remains
\code{not\_certified}. Phase 6 is the extreme case: its agent declares no tool and
runs a single turn, so that the report is assembled and the model writes only a
bounded analyst note.

\marginref{Phase modules: \code{src/agent/phases/registry.py}, \code{PHASE\_MODULES}. Phase-1 projection: \code{GraphPhase.\_build\_graph\_evidence\_projection}.}

\subsection{Graph and run context}

The graph is built in two layers: a loader reads the infrastructure description
into data objects, then populates a backend. The backend is an abstract interface
with a network-graph implementation, so that another engine can be substituted
without touching the tools. Two details matter for reading results: the neighbour
query is the sorted union of successors and predecessors, and path search uses an
undirected view so that a path is found in both directions.

The state belongs to the run, not to the process. Each pipeline creates an
immutable run context carrying the backend, the infrastructure, the vulnerability
reports, the risk scores, the scenario or discovery topology and the weighted
graph. The graph tools retrieve that context without changing the arguments
exposed to the model; thread workers capture the owner explicitly rather than
relying on implicit propagation.

\marginref{\code{src/loader.py} and \code{src/graph\_backend.py}; context in \code{src/agent/core/graph\_context.py}.}

\subsection{The tool governance boundary}

All profiles and all modes go through the same wrapper. On each call it checks the
previous piece of evidence, the stop request and the budgets; increments a counter
and refuses the call if the ceiling is reached; assigns a reference; applies the
scope refusals; executes; and archives an entry, including for a refusal or an
exception.

\begin{table}[H]
\centering
\caption{Controls applied to every tool call and the consequence of a failure.}
\label{tab:executor}
\small
\begin{tabular}{@{}p{3.6cm}p{5.2cm}p{5.0cm}@{}}
\toprule
Control & Mechanism & Consequence of a failure \\
\midrule
Evidence integrity & a flag set when the archive cannot be written & the run stops, and is not verifiable \\
User stop & a stop event, in the experimental scope & the run stops \\
Time budget & monotonic clock against the limit & the run stops \\
Monetary budget & the cost tracker, on observed usage & the run stops \\
Call budget & a counter taken under the archive lock & the run stops \\
Scope & the destination fields against the run's target network & the call is refused and archived \\
Free execution & the execution tool & refused as unverifiable scope \\
\bottomrule
\end{tabular}
\end{table}

Two properties matter for interpreting a result. A failure to write the log is
never ignored: it interrupts the run rather than leaving untraced actions, which
is the difference between an audit and a log that only records what went well. And
the cost limit applies to observed usage, not to a reservation of calls still in
flight, so a run can exceed its monetary budget by the cost of the calls it had
already started.

\marginref{Implementation: \code{wrap\_tool} and \code{check\_execution\_limits}, \code{src/agent/core/executor.py}.}

\subsection{Artefacts, cost tracking and retention}

The run folder is created exclusively under the output tree; in the event of a
collision a hexadecimal suffix is appended and no earlier folder is overwritten.
The run metadata records the commit, the policies and the limits, and the cost
summary is written in the finalisation block, hence even when the run raises. A
submission bundle is a zip equipped with a SHA-256 manifest that is re-checked
after writing; the download route excludes private files and hidden entries.

The price of a phase is resolved once and recorded with its source, so that a
completed run cannot be revalued after the fact; the summary distinguishes the
cumulative duration of the agents from the wall-clock time, which are not the same
measure.

\marginref{Collection contract: \code{src/benchmark/artifacts.py}. Pricing: \code{src/agent/cost\_tracker.py}.}

\subsection{What the architecture does not guarantee}

The isolation is per run for the deliverables and the graph, but not entirely so:
the phase services may still call the common engine, and the simultaneity of two
scenarios on a shared laboratory is not guaranteed. The laboratory reservation is
active only when its variable is set, and is reentrant only within one thread.
Without explicit parameters, the interactive pipeline imposes neither a call
ceiling nor an analysis timeout per device, so the budgets of
table~\ref{tab:executor} are conditional rather than structural.

Three points could not be established by reading the code, and are therefore not
asserted: no archiving or compaction mechanism for run folders was found, so
retention appears to be that of the file system; one evidence file is read by two
modules but its main producer was not identified; and the exact field content of
the provider diagnostics log was not verified line by line. Finally, a completed
status means that the runner finished its work, never that the target was
exhaustively covered nor that an access was proven.

\section{Audit protocols and evidence contracts}
\label{sec:protocoles}

\textbf{In brief.} Each family splits what it tests into three artefacts that the
code keeps apart. The \emph{briefing} states only the task, never the expected
answer. The \emph{private oracle} states which state was installed and which fact is
observable, and is written before the policy runs. The \emph{offline evaluator}
re-reads the log and the final decision, and recomputes the outcome. A success is
therefore never a declaration; it is a fact that can be recounted from a journal
the policy did not write.

\textbf{What this buys, and what it does not.} It buys freshness and attribution: a
witness drawn at each trial proves the value was actually read, and an observed
disclosure that contradicts a protected oracle invalidates the trial instead of
being counted as a success. It does not buy generalisation: each family covers one
protocol, and the V0 control of fidelity covers MQTT only.

\marginref{Contracts: \code{ssh\_reference.py}, \code{mqtt\_reference.py}, \code{api\_reference.py}, \code{v0\_control.py}, \code{mqtt\_formats.py} under \code{benchmarks/experiments/}.}

\subsection{Implemented families}

The repository contains five families, each with a centralised contract file, a
deterministic arm, a bounded LLM arm and a launcher reserved for the laboratory.

\begin{table}[H]
\centering
\caption{Families, what they measure and what they leave out of scope.}
\label{tab:familles}
\small
\begin{tabular}{@{}p{2.4cm}p{3.0cm}p{4.0cm}p{3.8cm}@{}}
\toprule
Family & Contract & What it measures & What it does not measure \\
\midrule
SSH & \code{ssh\_reference.py} & Whether a dictionary credential reads the witness object that the legitimate control also reads & Fixture reset, provenance of the model \\
MQTT & \code{mqtt\_reference.py} & Whether an anonymous subscriber receives the fresh private witness; in a healthy posture, the evaluated refusal & Purge of retained messages, broker reset \\
API & \code{api\_reference.py} & Coverage of the four authorisation boundaries & The baseline limit is \emph{declared}: the script reads no collection \\
V0 & \code{v0\_control.py} & Fidelity of execution: a public procedure executed and proved & An assumption about the model \\
H2 & \code{mqtt\_formats.py} & Reading of a frozen MQTT presentation & Network collection, field semantics, statistical independence \\
\bottomrule
\end{tabular}
\end{table}

The limit is not merely editorial: it is recorded in each result, under the key
\code{scope}, and in \code{measurement\_status: evidence\_only\_not\_admitted}.

\subsection{A three-part contract}

Each family separates three artefacts, and the separation is enforced by the
code.

\paragraph{The public briefing} contains only the task.
\code{ssh\_reference.validate\_briefing} requires exactly the expected set of
fields, which forbids adding a posture or a verdict. The SSH dictionary cannot
contain the control identity, and \code{mqtt\_reference.validate\_briefing}
refuses a public subject equal to the private subject. The posture is never
written in the briefing: it is inferred from the oracle.

\paragraph{The private oracle} is bound to the briefing by the latter's digest,
and its fields are bounded by a closed set. It states which state was installed
(F0 healthy or F1 vulnerable) and which fact is observable; it never states which
decision the policy should return.

\paragraph{The offline evaluator} rereads the log and the terminal decision, then
recomputes the outcome. The evidence file is \code{tool\_calls.jsonl} for SSH,
MQTT, API and V0; it is \code{journal.jsonl} for H2, reattached by
\code{Journal.replayed}, which refuses a record originating from another dossier.
Each record carries a reference produced by the shared executor
(\code{src/agent/core/executor.py}).

\subsection{The oracle: construction and guarantees}

The oracles are produced by \code{oracle\_for} functions, present in three
modules. SSH builds its own inline, from \code{case\_configuration}, which draws
the object with \code{secrets.token\_hex(32)}; MQTT produces two fresh witnesses
per trial; the API randomises organisations, users, objects and collections at
each trial.

Three barriers prevent deriving the answer from the policy. The first is
structural: the evaluator receives only the briefing, the oracle, the journal and
the decision, and the oracle is written before the arm runs. The second is
cross-validation: a briefing that would name an oracle file fails, and an unbound
oracle fails. The third is the contradiction verified over the entire journal:
under a protected oracle, an observed disclosure makes the trial
\code{oracle\_conflict} rather than a success. The oracle therefore cannot be
adjusted after the fact.

The witness (\code{canary} in the code) is a fresh value drawn at each fixture
start, whose digest is in the oracle and whose value appears nowhere in the
briefing. It proves freshness and qualifies the control: the legitimate control
must have read the value. Without a witness, an observed refusal might be mere
silence, which MQTT explicitly refuses.

\subsection{The offline evaluator: recount and statuses}

The terminal decision is never taken on trust alone. The evaluator requires that
its evidence references be a subset of the journal, and otherwise returns
\code{outcome: unknown}. It then recomputes per observation type: an archived
request must be identical to the requested request, and the payload is extracted
only if the return code and the framing are consistent. Unbound evidence does not
count, even when cited.

In F1, success requires a finding, a proven weak access and a functional control.
In F0, it requires the absence of a finding, a functional control and the explicit
refusal of each pair of the dictionary. The code distinguishes \emph{non-success}
(trial carried to completion, contract not satisfied) from \emph{missing} (evidence
or decision absent). The distinction is not cosmetic: a missing item does not
enter a success-rate denominator, a non-success fills it. In V0,
\code{instrument\_defect} is added when the script prevails and the LLM fails:
this is a defect of the instrument, not a weakness of the model.

\subsection{The V0 control}

V0 is the validity gate, not a hypothesis. The task is solvable without
inference: every subject, identifier and window appears in the briefing.

A \emph{saturating control} is a control that both arms pass everywhere: it
detects no defect. The manifest
\code{benchmarks/experiments/v0-control-manifest.json} freezes \code{X = 0.9},
24 steps and 3 repetitions, and declares \code{saturating: true} with its
reading: the control cannot detect a defect rarer than one failure in three. The
evidence is a run with three repetitions in which both arms succeed everywhere.
The threshold is read by \code{frozen\_threshold}, which refuses a manifest that is
not frozen or that was calibrated on a different number of steps.

\subsection{H2: three arms, one absorbed prediction}

H2 tests whether a model reads a presentation that no declared parser covers. The
eight dossiers of the corpus cross \code{\{reference, variant\}} and
\code{\{raw, normalized\}}, two dossiers per group. Three arms share the same
journal and the same verdict vocabulary.

\begin{table}[H]
\centering
\caption{Success rate per group for the three H2 arms (measurement of 22 October).}
\label{tab:h2}
\begin{tabular}{@{}lccc@{}}
\toprule
Group & Declared parser & By convention & LLM \\
\midrule
\code{reference/raw} & 1.0 & 1.0 & 1.0 \\
\code{reference/normalized} & 1.0 & 1.0 & 1.0 \\
\code{variant/raw} & 0.0 & \textbf{1.0} & \textbf{0.5} \\
\code{variant/normalized} & 1.0 & 1.0 & 1.0 \\
\bottomrule
\end{tabular}
\end{table}

The LLM gain holds for one dossier in two. But the conventional arm covers both
without a model, since the convention suffices to read the variant. The control
therefore publishes \code{convention\_gain\_absorbed: true}: what the LLM gains
over the declared parser, a script that tries also gains. The prediction remains
literally upheld, but it no longer distinguishes the model from the declaration.

\subsection{The strict-v3 contract and confirmed metrics}

The contract \code{src/benchmark/strict\_v3.py} contains no scoring code: it
normalises the contract per vulnerability and derives a conservative value for
earlier ground truths. The \code{confirmed.*} metric is built in the comparator by
prefixing the funnel stage names (\code{src/benchmark/funnel.py}):
\code{candidates}, \code{filtered}, \code{confirmed}. The \code{confirmed} stage
denotes the deduplicated findings whose evidence has been independently validated
against tool outputs. The code states it: a confirmation whose evidence is invalid
is not a true positive. The stage is unavailable if the exploitation file or the
provenance journal is missing; it is never reconstructed from the survivors.

\subsection{What these contracts do not prove}

\begin{enumerate}
\item \textbf{Nothing about the ground truth.} The oracle is supplied by the
operator; each family records this in its scope. A false oracle makes all verdicts
false.
\item \textbf{Nothing about admission.} Campaign admission requires the explicit
adoption of the strict-v3 metrics, two independent families, three repetitions
per posture and an attested resource parity: conditions not met by the
calibrations.
\item \textbf{Nothing about arm parity.} SSH is not an identical task: on the LLM
side, the model reads the result and judges equality with the control itself,
while the deterministic arm compares the payloads; the evaluator recomputes
off-band. The budget-parity condition covers only declared budgets, not the
hidden calls of the script arm.
\item \textbf{Nothing about a rare defect.} V0 is saturating: with three
repetitions, it detects only a repeated defect.
\item \textbf{Nothing about the real space.} H2 holds for eight hand-written
dossiers, a single variant schema and a single property.
\item \textbf{Nothing about semantics or injection}: these are H4 and H6, outside
the scope of these contracts.
\item \textbf{Nothing about comparative cost}: the deterministic arm consumes no
tokens.
\item \textbf{Certain guarantees belong to the written protocol}, not to a code
assertion: the arm count of H2 and the invariant that an observation cannot answer
two steps.
\end{enumerate}

\section{Harnesses, and a sandbox that measures nothing}
\label{sec:harnais-sandbox}

\textbf{In brief.} A harness is the instrument that puts two policies in front of
the same situation: one contract file per family, a deterministic arm, a model arm,
a controller for the laboratory and the evaluator. The same anatomy is repeated
for SSH, MQTT, the API and the V0 control. In October the life cycle of these
objects was unified into one construction function, so that a run missing part of
the governance surface is refused at construction rather than misbehaving later.

\textbf{The sandbox and its honest limit.} A sandbox was added that runs the
deterministic arms on any workstation and in continuous integration, without a
broker, a model or a laboratory. It reproduces the H2 lesson offline: the declared
parser fails where a script reading by convention succeeds. It \emph{instruments}
a contrast; it measures nothing about a model, and its guard observes Python events,
not system calls. The report therefore says \code{guard\_clean}, never
``isolated''.

\subsection{Anatomy of a harness}

A harness is a set of files with one responsibility each. For MQTT, which gives
its complete form: the contract validates the briefing, binds the tool to the
execution boundary, provides the deterministic arm and recounts the run through the
evaluator. The model arm is a separate file, and its H2 variant receives no parser.
The live controller builds the model, journals every transport call and writes the
decision, the cost summary and the run metadata. The deployment uses the
laboratory, submitting operations to the hypervisor.

The same anatomy is repeated for SSH, the API and the V0 control, with one
difference: the H2 arms --- the contract and its evaluator --- need neither a
deployment nor a broker, because their corpus is frozen. Its controller
\code{mqtt\_formats\_live.py} does require the laboratory runner and the model
catalogue, so it is not a sandbox instrument.

\subsection{The shared core}

The core exists for a precise reason, written in its header. The governance
boundary reads a fixed set of attributes from the run object with defaults; an
omission therefore manifests only when a tool call misbehaves. Each harness used to
rebuild that object by hand. The core replaces the copies with a single
construction, which enumerates thirteen required attributes --- the run directory,
three budgets, the archival lock, the call counter, the start clock, the limit reason, the
scope, the context, the policy and the cost tracking --- and refuses an incomplete
one.

One attribute deserves attention: the evidence-integrity flag. The limit check
reads it first and raises if it is set, and the archiver sets it when the archive
cannot be written. A harness that forgot it would pursue a run whose archive is
holed, and the result would look complete. The refusal is two-level: at
construction, and on any object through a replay check.

\subsection{The sandbox}

The sandbox proves, on a workstation as in continuous integration, that a harness
executes without a broker, without a model and without a laboratory: it loads the
frozen corpus of eight fixtures and, for each one, validates the fixture, opens a
journal, runs the chosen arm, writes the artefacts and recounts through the
existing evaluator. Five artefacts are produced per trial, and a report aggregates
them with a guard block.

The arm is an explicit parameter and an unknown name raises an error. This was a
correction: the first version read the arm from the corpus, which does not declare
it, so the conventional arm was unreachable and the report invented the policy it
believed it had tested.

\subsection{The guard, and what it does not see}

The guard relies on the audit hook of the interpreter. It stays installed but
enforces only inside an explicit context, and it is rebound at each run so that a
second run does not observe the first run's directory. It refuses sockets,
subprocesses, the process-spawning family, native library loading, and writes
outside the authorised directory.

Its honesty is itself a result worth keeping. A read-only audit broke the first
version with two measured escapes --- a spawned process and a native call through
the C library --- although the report announced a hermetic execution without
violation. Three changes followed: the key now says \code{guard\_clean}, which
states only that no listed event occurred; a guard block declares the mechanism,
the blocked events and the blind spots; and the isolation claim is explicitly
false. The last is false by construction, because an audit hook observes events of
the interpreter, not system calls.

\begin{table}[H]
\centering
\caption{What the guard covers, and what it leaves unobserved.}
\label{tab:guard-limits}
\small
\begin{tabular}{@{}p{6.2cm}p{7.6cm}@{}}
\toprule
Covered, with refusal raised & Unobserved \\
\midrule
Sockets: connection, resolution, sending, listening & Native code reached without the loading library, subprocess or spawn \\
Processes: subprocess, system, exec, fork, spawn & Writes through a descriptor opened before arming \\
Native libraries: loading, symbol resolution, call & System calls issued outside the interpreter \\
Rename and delete operations & Any behaviour after disarming: the hook remains installed but inert \\
Writes outside the authorised directory, on resolved paths & No guarantee of process isolation \\
\bottomrule
\end{tabular}
\end{table}

\subsection{What the sandbox demonstrates, and what it does not}

On the eight frozen fixtures, the declared-parser arm succeeds in six trials: the
two failures are the raw variants, where the declared parser does not cover the
schema and abstains, hence an erroneous conclusion. The arm reading by convention
succeeds eight times out of eight. Success is not counted on the conclusion alone:
the evaluator requires a correct conclusion, no contradiction and no erroneous
verdict.

The lesson is the same as H2. The failure does not stem from a lack of
intelligence but from a boundary of the declared parser, and a script that tries
restores the coverage. The sandbox instruments that contrast; it does not measure
it on a model. No model arm appears there, no pre-flight attestation is produced,
and no admissibility flag emerges from it.

Five limits must be read against those results. The sandbox involves no model, so
its success is a property of the harness, not an experimental result, and the
report carries \code{claims\_measurement: false}. The guard is a Python
observation, not a proof of isolation. The negative tests prove the refusals, but
only at unit level: no end-to-end test confirms that the installed hook is reached
during a run. The corpus is frozen, with one scenario and eight fixtures, which
suffices to instrument a contrast but not to estimate a rate. And the sandbox
bypasses the governance boundary it accompanies, calling the arm directly and
writing its own journal, so the core and the guard remain two instruments verified
in isolation.

\subsection{What the tests prove}

Twenty-four tests cover the core and the sandbox. What they prove:

\begin{itemize}
\item an incomplete run is refused at construction;
\item the call budget is raised by the wrapper rather than by the check;
\item the eight fixtures produce the five artefacts and a consistent report;
\item the arm is declared rather than read from the corpus;
\item the guard refuses writes outside the directory, the network, subprocesses
and three native-library events.
\end{itemize}

What they do not prove: no test executes a model arm, none measures an interval,
and none makes a sandboxed run attempt a forbidden operation, so the sanction is
proved at unit level rather than end to end.

\paragraph{Where this lives.}
\begin{itemize}
\item Harnesses, the shared core and the sandbox live under \code{benchmarks/experiments/}; tests in \code{tests/test\_harness\_core.py} and \code{tests/test\_sandbox.py}.
\item The core is \code{build\_run}; see also \code{REQUIRED\_ATTRIBUTES} and \code{require\_attributes}.
\end{itemize}

\section{Execution infrastructure and delivery pipeline}
\label{sec:infrastructure}

\subsection{Why a laboratory at all}
\label{sec:proxmox}

An audit cannot be measured on a real network. The result would depend on whatever
happens to be exposed that day, could not be repeated, and would not be
comparable between two runs. LANCE therefore audits networks that it has itself
built, from a description, and throws away afterwards.

A scenario is that description: which devices exist, which services they run,
which weaknesses are installed, and which connections are allowed. Deploying it
means materialising those devices as virtual machines on a hypervisor, so that the
audit sees a real network rather than a simulation, and destroying them afterwards
so that the next trial starts from the same state.

\textbf{Proxmox} provides that hypervisor. It is the host on which the scenario
devices are created as LXC containers --- lightweight, unprivileged Linux
containers, each acting as one device of the scenario: a web service, an IoT
service, a gateway, or any other role the scenario describes. The scenario never
runs on the workstation and never touches a real network: it exists only on the
Proxmox host, for the duration of a trial.

\textbf{Ansible} plays the other half. It configures each container from the
scenario description --- installs the service, sets the credentials, injects the
weakness --- and resets it to a clean state before the next trial. Deployment,
configuration and reset are the three operations that make a trial repeatable; the
hypervisor alone would give none of them.

The cycle is therefore: deploy the scenario as containers, configure and inject,
run the audit, destroy. A run that skipped the reset would inherit the previous
trial's state, and its result would no longer be attributable.

\begin{figure}[H]
\centering
\begin{tikzpicture}[
  box/.style={draw, rounded corners=2pt, align=center, font=\footnotesize,
              minimum height=10mm, text width=30mm, inner sep=3pt},
  arr/.style={-{Latex[length=1.6mm]}, font=\scriptsize, draw=black!70},
  node distance=7mm and 12mm,
]
\node[box] (scen) {Scenario\\[-1pt]\scriptsize devices, services,\\[-1pt]\scriptsize weaknesses, links};
\node[box, below=10mm of scen] (ans) {Ansible\\[-1pt]\scriptsize configure, inject, reset};
\node[box, below=10mm of ans] (px) {Proxmox host\\[-1pt]\scriptsize LXC containers:\\[-1pt]\scriptsize web, IoT, gateway};
\node[box, right=22mm of px] (agent) {Audit agent\\[-1pt]\scriptsize plans and runs\\[-1pt]\scriptsize the tests};
\node[box, above=10mm of agent] (eval) {Evaluator\\[-1pt]\scriptsize recomputes from\\[-1pt]\scriptsize the journal};
% flows
\draw[arr] (scen) -- node[left]{\scriptsize deploy} (ans);
\draw[arr] (ans) -- (px);
\draw[arr] (px) -- node[above]{\scriptsize tests, observations} (agent);
\draw[arr] (agent) -- node[right]{\scriptsize journal} (eval);
\draw[arr, dashed] (agent.west) to[bend left=35] node[below]{\scriptsize destroy} (px.east);
\end{tikzpicture}
\caption{The laboratory cycle. A scenario is deployed as virtual devices on the
Proxmox host, configured by Ansible, audited, then destroyed. The evaluator reads
the journal, never the devices.}
\label{fig:labo}
\end{figure}

The platform itself runs in one unprivileged container on that same Proxmox host.
It carries the service, the runner and the artefacts. Two interfaces are attached:
one on the management network, in DHCP, and one on the scenario network, with no
gateway, which confines the scans to the isolated segment. An output firewall lets
through only what the platform needs --- name resolution, HTTP, SSH, the hypervisor
API, the tunnel of the overlay network and the package cache --- and rejects
everything else.

One detail of documentation matters for reproducibility: the host name appears in
the project's own guides but not in the versioned inventory, which is a template.
The only constant readable in the repository is the container identifier, 200. The
actual configuration lives outside the repository, and the report says so rather
than pretending otherwise.

\marginref{Playbook: \code{benchmarks/ansible/playbooks/deploy\_master.yml} (Debian 13, 20~GB, 4~GB memory, 2 cores, nesting, on boot).}

\subsection{The three instances}

Three instances share the container, one per branch, each with its own environment,
database, configuration and artefacts: the stable instance on the main branch, and
one working instance per development branch. The deployment script fixes the slot,
the port, the directory and the service, and rejects any other name. Every
instance answers a status route that announces its slot, its branch, its commit and
whether the laboratory lock is active; after start-up, the script checks that the
announced commit is exactly the one that was validated.

The lock is a single file, taken exclusively before any mutation and never deleted.
The same value is passed to the three instances, which is what allows them to
coordinate on one physical laboratory. Waiting is cooperative, in memory and
without arrival ordering, so it is an etiquette between instances, not a security
isolation between users.

\marginref{Slots and ports: 8501 for \code{main}, 8502 for \code{dev/1}, 8503 for \code{dev/2}; status route \code{GET /api/environment}; lock at \code{/var/lib/lance/lab.lock}.}

\subsection{The runner and the delivery pipeline}

The GitHub Actions runner is installed in \code{/opt/actions-runner}, registered
with the label \code{nato-master} and run as a systemd service.

The pipeline reads as four links. A \code{push} to \code{main}, \code{dev/1} or
\code{dev/2}, or a merge request into these branches, triggers
\code{.github/workflows/benchmark-integrity.yml}. The first job runs
\code{.github/scripts/ci\_scope.py}, which classifies the modified files and
publishes the \code{full} flag. The second builds the image of the isolated worker
and checks the absence of ground truth. The deployment jobs then call
\code{.github/workflows/update-master.yml}, which runs the deployment script with
the slot, the validated SHA, the branch and the remote repository.

\begin{table}[H]
\centering
\caption{Triggers and paths through the pipeline.}
\label{tab:circuit}
\small
\begin{tabular}{@{}p{3.1cm}p{5.6cm}p{4.8cm}@{}}
\toprule
Event & Path & Deployment \\
\midrule
documentation \code{push} & classification to \code{full=false}: diff of the previous commit and whitespace check & none \\
code \code{push} & full tests, scenario and Ansible checks, sandbox, worker image & \code{main} to 8501, \code{dev/*} to 8502 or 8503 \\
Merge request & same path according to classification & none \\
Manual trigger & \code{full=true} forced & chosen slot, or checks only \\
Version tag & \code{full=true} forced & \code{main} only \\
\bottomrule
\end{tabular}
\end{table}

The light path covers the Markdown, LaTeX, PDF and image files under
\code{docs/}, \code{objectifs/} and \code{research/}, as well as a closed list of
presentation files. The classification keeps deletions and the two paths of a
rename; any undetermined event falls back to the full path, and a failure blocks
the deployment. A documentation push therefore restarts no instance.

On the machine side, the update script fetches the SHA and applies a fast-forward
merge for \code{main}, or a detachment after checking for a clean tree for the
development branches; it then verifies that \code{HEAD} equals the validated SHA.
Before that, the service is stopped, the lock acquired, and the instance queried
in order to refuse mutating active or pending work. In the event of a failure after
stopping, the service remains stopped: no automatic rollback to an earlier version
is performed.

\subsection{Corrections to the deployment script}

The history of the deployment script is a short sequence of corrections made on
29 October 2026, the day the container was rebuilt and every automatic deployment
stopped working. Four families of correction came out of it, and each answers a
concrete failure rather than a matter of style.

First, the script assumed that it inherited the host's identity, the root user and
the presence of the neighbouring instances. A rebuilt container breaks all three
assumptions, so each became an explicit check. Second, the host identity is now
read from the overlay network interface rather than from a constant, which removes
the last value that a rebuild can invalidate. Third, privilege escalation loses the
environment: the path must be absolute, the script is re-executed through a shell,
and both the runner name and the deployment settings are carried explicitly.
Fourth, a download failure used to stop the deployment and leave the instance
stopped; it is now tolerated and named.

The same identity logic now serves the benchmark controllers, and a laboratory lock
variable allows an operator to suspend the check explicitly rather than by
disabling code.

\marginref{Commits: \code{6b0a9ff}, \code{433e2a3}, \code{a568c68} to \code{dcd5903}, \code{150eb8a}. Script: \code{scripts/deployment/deploy-instance.sh}.}

\subsection{The UMONS model and its relay}

The \code{ollama-umons} provider is reached through a local SSH transfer on the
Proxmox host, towards the university cluster (\code{docs/guides/umons-llm.md}). The
path was identified and restored on 4~October 2026, with a successful minimal
generation; the authentication observed is interactive. The restoration command is
documented, as are the listening checks. The guide is explicit: neither start-up
at boot, nor reconnection after the process ends; the administration of the
downstream GPU server remains to be confirmed with the administrator. The
fragility is established in the facts: the relay went down again on 22 October,
and it was the provider pre-flight that named the failure without executing any
arm.

\subsection{The outage of 29 October}

The outages actually traced in the repository for that day are those of the
container rebuild, which rendered obsolete the fixed address and the assumption of
a privileged runner. By contrast, \emph{no mention of a multipath default route on
the laboratory router exists in the repository}: the search on the corresponding
terms returns nothing. This fact, observed in session but not recorded, must not be
presented as documented. It constitutes a gap in the project's memory, since the
same outage will recur at the next router restart.

\subsection{Known fragilities}

\begin{itemize}
\item \textbf{Unreplicated failure points}: single Proxmox host, single container,
single runner. The runner registration tolerating failure, an absent runner does
not block the playbook but blocks every deployment.
\item \textbf{UMONS relay}: no supervised service, interactive authentication,
address still fixed in the examples of the guide although commit
\code{e54da24} claimed to have stopped fixing addresses in the documentation.
Restarting LANCE does not repair the relay.
\item \textbf{Cooperative lock}: without arrival ordering, in memory, not covered
for an Ansible script launched directly, an old checkout or a manual
intervention. This is not security isolation between collaborators.
\item \textbf{Manual recovery}: a failure after stopping leaves the service
stopped, with no automatic rollback.
\item \textbf{Host identity}: without a declared variable, the identity depends on
the presence of any address on the Tailscale interface; the script does not verify
that this host is indeed the one that should host the instance.
\item \textbf{Partial documentation}: the name of the Proxmox host and the
physical segment correspond only imperfectly to the playbook; the versioned
inventory is a template.
\end{itemize}

\section{The benchmark, and the distance left to a result}
\label{sec:benchmark}

\textbf{In brief.} The benchmark is a register of hypotheses, not a list of
scenarios. Six hypotheses are planned (H1 to H6), each behind a validity gate V0
that asks whether the instrument can reproduce a known control before anything
downstream is read. Measured to date: V0 passes on MQTT and saturates; H2 was
measured on 22 October and its prediction held, but the gain of the model was
absorbed by a script reading by convention; one family (C2) was rejected on paper,
and one (the vendor-register family) was built and rejected in an hour on
30 October.

\textbf{What this means for the question of the project.} The corpus was built
around what signature rules find: at most 2.1\,\% of its entries need
interpretation, and 84.4\,\% are detectable by signature. A duel on this corpus
therefore measures above all the quality of the rules, which is exactly what H2
observed. The conclusion is not that the model is inferior; it is that no situation
in which the difference could appear has been produced yet.

\subsection{The register, and the gate that is not applied}

Before every hypothesis the register places a fidelity control: if a task solvable
without inference fails on one of the two arms, no downstream result is readable.
The threshold is frozen at 0.9 with 24 steps, 3 repetitions and a saturating
control. The calibration pass returns a valid control on both arms, three times out
of three, and the register draws from it a \emph{partially passed} gate: the
instrument is readable for parity of execution and for a repeated defect, not for a
rare one.

Two caveats must be stated. The control covers MQTT only: the SSH harness keeps a
fixed credential dictionary and is therefore not an instance held apart. And the
threshold is not applied by the gate: the gate function reads the four outcomes on
``all trials succeeded'', while the threshold is computed separately and conditions
nothing. With three repetitions the two formulations coincide, but by numerical
coincidence, not by application of the threshold. This is one of the four
contradictions the conclusion confronts.

\begin{table}[H]
\centering
\caption{Status of the register on 2 November 2026.}
\label{tab:registre}
\small
\begin{tabular}{@{}p{1.1cm}p{3.4cm}p{8.2cm}@{}}
\toprule
Item & Status & What this means \\
\midrule
V0 & partially passed & valid control on MQTT; threshold frozen but not applied by the gate; SSH not covered \\
H1 & conditional, not opened & the identity-chaining family does not exist in the coverage plan \\
H2 & measured on 22 October & prediction held, gain absorbed by an arm by convention \\
H3 & not implemented & the public-contract families are planned \\
H4 & not opened & the vendor-register family was built and rejected on 30 October \\
H5 & conditional, planned & budget counter-test to be written \\
H6 & planned as a sub-test & graftable onto H2 or H3, without its own simulator \\
C2 & rejected & two adversarial reviews, dead on paper, no run \\
\bottomrule
\end{tabular}
\end{table}

\subsection{The corpus is built around what rules find}

The catalogue enumerates 29 scenarios; the derived index counts 288 expected
entries, of which 282 are classified as detectable by signature and 6 as belonging
to logic, that is 2.1\,\%. Two counts coexist and do not measure the same thing: the
\emph{mechanical} reading places every category except two in signature, while the
\emph{interpretive} reading judges, category by category, whether a well-designed
rule can conclude on its own, and finds 243 entries out of 288, that is 84.4\,\%.
Both are legitimate and must never be cited one for the other.

\begin{table}[H]
\centering
\caption{Distribution of the 288 entries by category and interpretive reading.}
\label{tab:corpus}
\small
\begin{tabular}{@{}lrl@{}}
\toprule
Category & Count & Interpretive reading \\
\midrule
\code{misconfiguration} & 70 & signature \\
\code{data\_exposure} & 54 & signature \\
\code{no\_authentication} & 50 & signature \\
\code{default\_credentials} & 41 & signature \\
\code{info\_disclosure} & 16 & signature \\
\code{code\_injection} & 13 & mixed \\
\code{insecure\_update} & 12 & mixed \\
\code{weak\_crypto} & 10 & mixed \\
\code{cve} & 9 & signature \\
\code{privilege\_escalation} & 4 & chaining \\
\code{broken\_access\_control} & 4 & logic \\
\code{missing\_header} & 3 & signature \\
\code{auth\_bypass} & 2 & logic \\
\midrule
Total & 288 & 243 signature, 35 mixed, 4 chaining, 6 logic \\
\bottomrule
\end{tabular}
\end{table}

The logic stratum occupies only three scenarios, and only 55 entries are attached
to a chaining of at least two hops. A global duel on this corpus would therefore
measure above all the quality of the rules, which is exactly what H2 observed.
Both counts are recomputed from the derived index, and the word \emph{signature}
means 282 entries mechanically and 243 interpretively.

\subsection{H2: a prediction that holds and demonstrates nothing}

The corpus contains eight fixtures, two per group, for 25 messages in total. The
measurement of 22 October 2026 used two arms and one repetition; the reference run
that followed used three arms per fixture --- 24 arms per repetition, 72 for three
repetitions.

On the raw variants, the model arm obtains 0.5 (one fixture won out of two), the
declared parser 0.0, and the third arm, which reads by convention without any key
from the corpus, covers 1.0. The prediction holds in the literal sense: the model
does better than a declared parser that cannot read a schema it does not know. But
its gain is absorbed by a convention that no one had to declare, and a convention
is a rule.

\subsection{Inference: repaired, and still idle}

The inference tool carried a structural defect. It computed a two-sided p-value
around zero, so every outcome returned the same figure: an almost constant series,
a highly dispersed series of the same mean and a series entirely negative all
produced a null p-value. A model performing strictly worse than the rules looked as
certain as one performing better.

\textbf{The defect was repaired on 30 October 2026.} The module is
\code{src/benchmark/inference.py}; the test is now one-sided
against the pre-registered threshold: an exact sign test up to seven families, a
paired bootstrap beyond, both centred on the threshold rather than on zero. Below
three families no p-value and no interval are emitted; a decision requires five
families, because five agreeing families are the first that can reach the usual
threshold at all.

\textbf{Repairing the computation did not make a conclusion possible.} A second
obstacle stands behind the first, and it is not a defect of the tool: the number of
independent families available. The trial configuration includes commit,
fingerprints and budgets, so two distinct families form two distinct configurations;
each configuration contains one family, and the inference requires several. No
figure can be produced from families that do not exist.

\subsection{A family built, then rejected in one hour}

On 30 October 2026 a family was built far enough to be measured, and rejected. The
record matters more than the rejection: it is the first time the project's own
protocol produced a decision instead of a plan.

\textbf{What was built.} Not the family: a probe. One fixture of three registers ---
one that commands something, one inert, one cache --- with real neighbour reads, test
writes, and every read attributed to a single write. The probe applies two arms: a
convention arm reading the unit and the change, and an arm reading the neighbours.
Neither an oracle, nor the closure test, nor the eight fixtures were built: the
rejection came first.

What existed before the probe was the design, and it resisted six recorded reviews.
Its addressing model was established against raw protocol frames rather than through
a client object, and reproduced by three readers. The earlier drafts had been
checked by reading; the frames checked the offset that the readings had assumed.

\textbf{The trap.} One register commands something and changes when written; one is
inert and does not change; the third is a cache --- it changes when written but
commands nothing, and propagates its value into its neighbours. The convention arm
classes the cache as a projection and scores two out of three. The arm reading the
neighbours scores three out of three and states why the cache is not a projection.

\textbf{Why it was rejected.} Reading the neighbours is a rule, not an
interpretation. Whether one places that rule inside the competent automation or
outside it, the family measures nothing. Inside, the automation succeeds and the
closure test rejects the family. Outside, the family measures an omission of the
automation rather than a limit of interpretation. There is no third reading.

The filter that would have caught it applies in minutes and costs nothing:

\begin{quote}
If the rule that solves the case can be written in one line, a script will have
it. Do not build the family.
\end{quote}

\textbf{What survives the rejection.} Three findings remain, and they are what the
next attempt inherits. First, absence can be deduced: a register that does not
change when written, after a declared test, is attested as non-commanding without
any observation labelled ``absent'' in the journal --- labelling would have put the
answer in the journal. Second, a journal must be attributable: no observation may
belong to two causes, and the probe found three defects of that kind. Third, the
protocol works: six reviews validated a design that one hour of observation
rejected, and the observation cost less than a campaign.

\textbf{What the rejection does not prove.} It does not refute H4. It shows that
this mechanism, on this corpus, is crossed by a generic rule. A different mechanism
could survive, and any future family must demonstrate, before construction, that no
generic rule crosses it.

\subsection{Campaign admission}

The admission module combines pair conditions and campaign conditions. The
condition for explicit adoption of the confirmation metrics is a parameter that
defaults to false and is never set to true by the summary nor by the command line.
The justification is a contract argument: the driver harnesses do not emit these
counters. This is accurate about the current state, but not a proof of
unsatisfiability: the condition would become true again the day a harness publishes
these counters. Today, admission is therefore refused because that condition is not
satisfied.

A second point, verified during drafting, concerns budget parity. The condition
requires that the first value exists and that the two are equal, so a missing budget
is correctly refused. The module does not, however, verify in its other conditions
that a missing budget is distinguished from a non-comparable one.

\subsection{What is publishable today}

\textbf{The publishable element is the method, not a result of superiority.} The
repository supports no claim of comparative performance. What holds is consistent
and verifiable:

\begin{itemize}
\item a falsifiable register of hypotheses, with the contrast declared before the
fixture, a private oracle and a measurement recomputed from the log;
\item a measurement lesson from H2: a prediction can hold literally and demonstrate
nothing, because an arm by convention absorbs the gain;
\item a clean negative result on C2, documented before any grey area;
\item a defensible negative result on the vendor-register family, obtained in one
hour instead of a campaign;
\item an honest inventory of the holes: a reduced logic stratum, an unknown protocol
absent, noise not encoded.
\end{itemize}

\subsection{What is not demonstrated}

Must not be written as established: the superiority of the model; a significant
gap, since no independent family exists to compute one; generality, H2 covering one
schema, one question and one model; a cost gain, the exploratory pilot counting one
pair without an interval; the validity of the threshold, which the gate does not
apply; and the integrity of the corpus, the official healthy control not being in
the catalogue.

The other active studies are recorded as qualifying the instrument rather than
showing an advantage: the plan carries 44 records and several hundred conditional
arms, of which the qualifications actually executed --- SSH in October, MQTT a few
days later --- validate the instrument. The three-repetition pilot yields twelve
admissible pairs and computes no interval. The formulation to keep is the one of the
historical study on agent versus automation: a real descriptive result, no general
superiority established.



% ---------------------------------------------------------------------------
\section*{References}
\addcontentsline{toc}{section}{References}

Internal sources are cited by path, with the date of the state they describe. No
external publication is claimed for a result that does not exist.

\begin{itemize}
\item LANCE project. \emph{Architecture, protocols and current state, report as of
30 October 2026}. \href{../../research/2026-10-30-architecture-protocoles-avancement/README.md}{\nolinkurl{research/2026-10-30-architecture-protocoles-avancement/}}.
\item Rejection of the pilot family \texttt{s\_registres\_constructeur}, 30 October
2026. \href{../../research/2026-10-30-famille-pilote-registres-constructeur/compte-rendu-rejet.md}{\nolinkurl{research/2026-10-30-famille-pilote-registres-constructeur/}}.
\item Hypothesis register and the V0 gate.
\href{../../research/2026-10-21-hypotheses-benchmark/README.md}{\nolinkurl{research/2026-10-21-hypotheses-benchmark/}}.
\item H2 measurement protocol and runs.
\href{../../research/2026-10-21-hypotheses-benchmark/h2-protocole.md}{\nolinkurl{research/2026-10-21-hypotheses-benchmark/h2-protocole.md}}.
\item C2 rejection.
\href{../../research/2026-10-21-c2-couverture/README.md}{\nolinkurl{research/2026-10-21-c2-couverture/}}.
\item Derived corpus index (\code{benchmarks/labels/derived\_index.json}) and its
generator, \code{src.benchmark.derived\_index}.
\item V0 fidelity control manifest (\code{benchmarks/experiments/v0-control-manifest.json}).
\item Reference material for the current behaviour: \code{docs/architecture/run-modes.md},
\code{docs/benchmark/derived-index.md}, \code{docs/guides/development-instances.md}.
\end{itemize}

\section{Conclusion}
\subsection{What this report can say, and what it cannot}

The project assembles three instruments that do not offer the same guarantees, and the main source of confusion is to treat them as interchangeable. Section~\ref{sec:etat} gave the state as a whole; this conclusion states what follows from it.

\begin{table}[H]
\centering
\caption{What each instrument guarantees and does not guarantee.}
\label{tab:instruments}
\small
\begin{tabular}{@{}p{3.4cm}p{4.6cm}p{5.0cm}@{}}
\toprule
Instrument & Guarantees & Does not guarantee \\
\midrule
Execution boundary (\code{executor.py}) & Each tool call is logged, bounded and compliant with the scope & System-level isolation; the guards are application-level \\
Sandbox (\code{sandbox.py}) & No listed Python event occurs; the deterministic arms run outside the laboratory & Native isolation; any model measurement \\
Calibration harness & Contract, separate oracle, evaluator that re-reads the log & Campaign admission; exact parity of the arms \\
\bottomrule
\end{tabular}
\end{table}

The first instrument makes the audit traceable. The second makes the harness portable. The third makes the comparison possible. None of them proves that a model does better than a script, which was the object of the programme.

\subsection{The contradictions identified between sections}

Comparing the parts brought out four discrepancies, each verified against the code before being retained here.

\paragraph{The threshold of the V0 check is not applied by the gate.} The protocol framework asserts that the gate applies \code{X = 0,9}; the benchmark part verifies that the gate is computed on ``all successes'' and that the threshold check is independent. The two formulations coincide at three repetitions but diverge beyond that. The report retains the reading of the code.

\paragraph{The H2 corpus does not encode the described failure mode.} The protocol states that, for a given record, abstention is the only correct answer; reading the oracle shows that it expects an exposure on the first message and a non-exposure on the third. The by-convention arm moreover converges exactly on the oracle. The protocol therefore describes an intention that the corpus does not encode.

\paragraph{The blocking of inference is not the one commonly alleged.} The admission condition is cited as the cause; reading shows that it is satisfiable in principle, and that the real constraint is the absence of independent families. This shift matters: it changes the priority of the work.

\paragraph{The routing failure of 29 October is not recorded.} It accounts for six hours of interruption and the rebuilding of the container, but no document in the repository describes it. This is a gap in the project's memory, not a disagreement between sections.

\subsection{A common pattern: absorbed gain}

Three episodes follow the same shape. In the H2 family, the model's gain over the declared parser is taken up by a script arm that reads by convention. In the sandbox, the same contrast is reproduced offline, without a model. In the rejection of C2, the targeted flaw turns out to be expressible from public inputs, and therefore closeable by a competent automation.

The pattern can be stated thus: \emph{what distinguishes the model from the script, in the constructed families, is a limitation of the declaration rather than a capability of the model}. This is not a demonstration of inferiority; it is the absence of a situation in which the difference appears. Yet such a situation is precisely the object of the programme.

\subsection{What the project can assert today}

Four assertions are defensible, each with its limitation.

An audit produces verifiable traces: each tool call is archived, referenced and bounded, and a lost piece of evidence interrupts the run rather than leaving a silent gap. The limitation is that these are application-level guards.

A harness can be executed outside the laboratory, without a broker or a model, under a safeguard whose blind spots are declared. The limitation is that it measures nothing: it instruments.

A third script arm absorbs a gain attributed to the model on the most documented family. The limitation is that this holds for eight fixtures and one property.

The corpus is mostly detectable by signature, which makes a null difference expected by construction. The limitation is that this concerns the current corpus, not the targeted situations.

\subsection{What the project cannot assert}

No superiority of the model is established. No confidence interval can be computed. No generalisation is possible from the existing measurements. The inference module, which produced the appearance of a measurement it did not provide, was repaired on 30 October 2026: it now refuses to emit a figure it cannot defend. What remains is the obstacle behind it — the families that would give it something to compute do not exist.

\subsection{The three decisions that govern what follows}

A first decision concerns the independent families, now that the module can no longer produce an indefensible figure: to build them, applying the filter that one hour of observation established, or to publish a bounded methodological result and declare it as such. The filter matters more than the count: a family crossed by a generic rule measures nothing, and five of them measure nothing either.

A second decision concerns the opening condition: to write a competent reference automation, or to explicitly renounce the hypotheses that depend on it. This decision makes several families that are currently undecidable legible.

A third decision concerns the comparable arm. A comparison of the two policies supposes that both receive the same observations, and four mechanisms currently favour the rules: the phase-four probes only they may emit, a supplement that rewrites the scan files, the scanner findings added to every aggregation, and a benchmark score reserved for the full pipeline. Either give the model an equivalent access, or state that the comparison measures verification rather than interpretation.

These decisions together determine whether an inference can one day exist.

None of these decisions is technical in the narrow sense. All of them concern what one accepts not to be able to say.

\subsection{A legible conclusion}

The system is built, instrumented and documented. Its chain of evidence is sound: a separate oracle, an evaluator that recounts, a log that does not forgive a lost piece of evidence, a safeguard that declares its blind spots. The defects identified in this report were found by reading what the code does, not what it claims, and that is the best evidence that the verification discipline works.

The scientific question, however, remains open. The programme asked when a model does better than a script with identical tools; the available data do not answer it, and the current arrangement would not allow an answer even with more machine time. This observation is not a failure of the work accomplished: it is the result of having looked at things squarely.

The most useful continuation is not to add a family, but to make falsifiable what today claims to be so.


\end{document}
